% CLAUDE REPORT::WRONGPROOF-BRIDGE-0001::2026-11-13
% !TEX program = xelatex
% Compiles with XeLaTeX or pdfLaTeX (ASCII source).
\documentclass[11pt]{article}
\usepackage[margin=1in]{geometry}
\usepackage{amsmath,amssymb,amsthm}
\usepackage{xurl}
\newtheorem{theorem}{Theorem}
\newtheorem{proposition}[theorem]{Proposition}
\newtheorem{corollary}[theorem]{Corollary}
\newtheorem{lemma}[theorem]{Lemma}
\theoremstyle{definition}
\newtheorem{definition}[theorem]{Definition}
\theoremstyle{remark}
\newtheorem*{remark}{Remark}
\newcommand{\SAT}{\mathrm{SAT}}
\newcommand{\EF}{\mathrm{EF}}
\newcommand{\FP}{\mathrm{FP}}
\newcommand{\WP}{\mathrm{WrongProof}}
\newcommand{\Refl}{\mathrm{Refl}}
\newcommand{\status}[1]{\textbf{[#1]}}

\title{CLAUDE WRONGPROOF-BRIDGE-0001\\[2pt]
\large Does the 2026 provable-reductions theorem give a weaker bridge\\
$\SAT\in\mathrm P\Rightarrow \overline W_R(N)\le\mathrm{poly}(N)$?\\
An exact composition, and why it adds no leverage}
\author{CLAUDE (scheduled routine, independent of the LSPV-1..5 roles)}
\date{2026-11-13}

\begin{document}
\maketitle

\noindent\textbf{Status.} P versus NP remains \textbf{OPEN}. This is an additive report written by a
scheduled Claude routine. It is not an LSPV run. It changes no schedules, allocations or ownership, and it does not edit
root \texttt{main.tex}, because \texttt{AGENTS.md} reserves the ownership regions for LSPV-1..5.
LSPV-1 (bridge lead) and LSPV-5 (red team / synthesis) may adopt or reject it.

\paragraph{One-paragraph summary.}
The shared objective is the worst-case minimum derivation cost
$\overline W_R(N)=\max_{F\in\mathrm{UNSAT},\,|\langle F\rangle|\le N}\min_{\pi:F\vdash_R\bot}\operatorname{cost}_R(\pi)$,
together with the weakest sufficient bridge
$B(R):\ \SAT\in\mathrm P\Rightarrow\overline W_R(N)\le\mathrm{poly}(N)$.
The literature preflight found one new clue relevant to $B(R)$: Fleming--Grosser--Pitassi--Robere,
\emph{Provable Reductions in TFNP} (ECCC TR26-109, June 2026; cited below as FGPR).
The corpus lists it only as a bibliography item (LSPV-1/361); no report applies its theorems.
Composing FGPR Theorem~5.2 and Proposition~3.8 with the solver-induced proof system of a
hypothetical polynomial-time SAT decider $A$ gives a sufficient condition for $B(\EF)$, and hence
for the project's $R_{\mathrm{mix341}}\ge_p\mathrm{ER}$. The condition is that an explicit search problem,
which is in $\FP$ unconditionally, be \emph{$\EF$-provably} in $\FP$.
\textbf{Main finding (PROVED/EXACT, given the cited theorems):} under $\SAT\in\mathrm P$ this condition is
\emph{equivalent} to the Krajicek--Pudlak reflection condition that LSPV-1/361 and the Codex
2026-11-11 audit already record. Both are equivalent to $\EF$ being p-bounded.
So the FGPR route is an exact reformulation. It does not give a weaker bridge.
The only real change is a sharper diagnosis: $B(\EF)$ can fail only if
\emph{every} polynomial-time SAT decider is correct but not $\EF$-provably correct, which is the same kind of algorithm as AKS.
The load-bearing arrows, namely an unconditional proof of $B(R)$ and a superpolynomial lower bound
for the same fixed $R$, remain \textbf{OPEN}.

\section{Objects and conventions}

\begin{definition}[Solver-induced system]
Let $A$ be a deterministic polynomial-time algorithm that correctly decides SAT for CNF formulas,
and suppose one exists. Define the Cook--Reckhow system $Q_A$ for UNSAT CNFs as follows. A $Q_A$-proof of $F$ is either
$(0,\Pi)$, where $\Pi$ is an $\EF$-refutation of $F$, or the single symbol $(1)$.
The verifier accepts $(1)$ iff $F$ is a well-formed CNF and $A(F)=\textsf{UNSAT}$.
\end{definition}
$Q_A$ is sound and complete because $A$ is correct. Its verifier runs in polynomial time. Trivially $Q_A\ge_p\EF$, and every UNSAT
$F$ has a $Q_A$-proof of size $O(1)$, so $Q_A$ is p-bounded \status{PROVED/EXACT}.
This is the system $S_A$ of the Codex 2026-11-11 audit with an $\EF$ component added.
For a system $P$, write $\Refl_P$ for its reflection principle (FGPR Def.~2.4) and
$\WP_P$ for the false-clause search problem of $\neg\Refl_P$ (FGPR Def.~5.1). Given $(F,\alpha,\Pi)$,
the task is to output a clause of $F$ that $\alpha$ falsifies, or an error in the claimed $P$-proof $\Pi$.
$\langle\EF,R\rangle$ is the FGPR system (Def.~3.5): a refutation of $F$ is a mapping reduction from
$\mathrm{Search}_F$ to $R$, together with an $\EF$ proof of its correctness.
$R$ is \emph{$\EF$-provably in $\FP$} (FGPR Def.~3.7) if, for each input length, a polynomial-size circuit
solving $R$ comes with polynomial-size $\EF$ proofs of its correctness.

\paragraph{Imported theorems \status{KNOWN THEOREM}.}
\begin{itemize}
\item[(K1)] FGPR Thm.~5.2: if $P\ge_p\EF$ has polynomial-size proofs of $\Refl_P$, then
$P\equiv_p\langle\EF,\WP_P\rangle$.
\item[(K2)] FGPR Prop.~3.8: if $R\in\mathrm{TFNP}$ is $\EF$-provably in $\FP$, then $\langle\EF,R\rangle\equiv_p\EF$.
\item[(K3)] FGPR Sec.~1 remark: $\WP_P\in\FP$ for every Cook--Reckhow $P$. The algorithm checks whether $\alpha$
falsifies a clause, and otherwise locates an invalid step of $\Pi$ by running the verifier.
\item[(K4)] Krajicek--Pudlak (1990): if $\Refl_P$ has polynomial-size $\EF$ proofs, then $\EF\ge_p P$.
$\EF$ has polynomial-size proofs of $\Refl_{\EF}$.
\item[(K5)] Cook--Reckhow (1979): some propositional proof system is p-bounded iff $\mathrm{NP}=\mathrm{coNP}$.
\end{itemize}
FGPR state (K1) and (K2) for tautologies. The project measures refutations of CNFs.
The standard Tseitin translation is $\EF$-polynomial in both directions, so the two conventions are interchangeable here.
This is the only convention step in the argument and is marked as such.

\section{The composition}

\begin{theorem}[WrongProof form of the bridge]\label{thm:wp}
\status{PROVED/EXACT relative to (K1)--(K4) and the Tseitin convention.}
Suppose $\SAT\in\mathrm P$ via a correct decider $A$. Then the following are equivalent:
\begin{enumerate}
\item[(a)] $\WP_{Q_A}$ is $\EF$-provably in $\FP$;
\item[(b)] $\Refl_{Q_A}$, equivalently $\Refl_{S_A}$, has polynomial-size $\EF$ proofs;
\item[(c)] $\EF$ is p-bounded, that is, $\overline W_{\EF}(N)\le\mathrm{poly}(N)$ with size as cost;
\item[(d)] (a) holds for \emph{every} correct polynomial-time SAT decider $A'$ in place of $A$.
\end{enumerate}
\end{theorem}
\begin{proof}
(a)$\Rightarrow$(c). $Q_A$ is p-bounded. Its reflection formulas are tautologies of size
$\mathrm{poly}(n,m,s)$, because its verifier is a polynomial-size circuit, so each has an $O(1)$-size $Q_A$-proof (tag $1$).
Hence (K1) applies: $Q_A\equiv_p\langle\EF,\WP_{Q_A}\rangle$. By (a) and (K2) this equals $\EF$ up to $\equiv_p$.
So $\EF$ simulates a p-bounded system and is itself p-bounded.

(b)$\Rightarrow$(c). By (K4), $\EF\ge_p Q_A$, which is p-bounded.

(c)$\Rightarrow$(a),(b),(d). If $\EF$ is p-bounded, every polynomial-size tautology family has
polynomial-size $\EF$ proofs. Take the correctness formulas of the (K3) circuit for $\WP_{Q_{A'}}$, for any correct $A'$.
They are true because $A'$ is correct, and they have polynomial size, which gives (a) and (d). Similarly
$\Refl_{Q_A}$ is a polynomial-size tautology family, which gives (b). $\Refl_{S_A}$ is the tag-$1$ half of $\Refl_{Q_A}$;
the $\EF$ half is (K4). (d)$\Rightarrow$(a) is immediate.
\end{proof}

\begin{corollary}[Bridge reformulation for the project's calculus]
\status{CONDITIONAL on the corpus fact that $R_{\mathrm{mix341}}$ contains paid ER with polynomial cost overhead.}
\[
B(R_{\mathrm{mix341}})\ \Longleftarrow\ B(\EF)\ \Longleftrightarrow\
\Bigl[\SAT\in\mathrm P\ \Rightarrow\ \exists\text{ correct poly-time }A:\ \WP_{Q_A}\text{ is }\EF\text{-provably in }\FP\Bigr].
\]
\end{corollary}
\begin{proof}
The equivalence follows from Theorem~\ref{thm:wp}, (a)$\Leftrightarrow$(c), applied under the hypothesis.
ER and $\EF$ are p-equivalent \status{KNOWN THEOREM}. Unit-node cost is at most proof size, and LSPV-1/360 F3 records ER as a paid
subsystem of $R_{\mathrm{mix341}}$. Together these give $\overline W_{R_{\mathrm{mix341}}}\le\mathrm{poly}(\overline W_{\EF})$.
\end{proof}

\paragraph{Classification \status{FAILED TRANSFER as leverage; EXACT OBJECT TRANSFER as theorem}.}
The FGPR condition is exactly the reflection condition (b) already recorded in LSPV-1/361 (eq.~(29)
of the Codex 2026-11-11 transfer report). The search problem $\WP_{Q_A}$ is in $\FP$ for free, by (K3),
but the bridge needs that membership to be \emph{$\EF$-provable}, and that requirement is the whole difficulty.
Item (d) is a mild new observation: under $\SAT\in\mathrm P$ it makes no difference \emph{which} decider is used.
"Free knowledge of the best algorithm" is therefore as harmless on the bridge side as "free knowledge of the best derivation" is on the
lower-bound side. It also gives no help.

\paragraph{Sharpened diagnosis \status{PROVED/EXACT, contrapositive of Thm.~\ref{thm:wp}}.}
If $\SAT\in\mathrm P$ and $B(\EF)$ fails, then every correct polynomial-time SAT decider $A$ is
\emph{$\EF$-unprovably correct}, in the sense that $\WP_{Q_A}$ is not $\EF$-provably in $\FP$. FGPR note that the only known
natural polynomial-time algorithm whose correctness is believed unprovable in $S^1_2$ is AKS primality testing.
An unconditional proof of $B(\EF)$ is therefore exactly a proof that a hypothetical polynomial-time SAT algorithm can always be
replaced by an $\EF$-provably correct one: a \emph{program-to-provable-program normalization} for SAT.
No such normalization theorem is known for $\FP$ problems in general. FGPR's own examples, where $R\in\FP$ but $\EF$-provable membership would make
$\EF\ge_p G_1$, show that it fails for some FP problems unless proof-complexity collapses occur.
This does not refute it for SAT.

\section{Required audit sections}

\paragraph{CORPUS REVIEWED.}
\texttt{AGENTS.md} (all directives, including the 2026-11-09 knot/link handoff, read but not acted on because this is not an LSPV run);
root tracker heads \texttt{proof2.txt} (status OPEN) and \texttt{proof1.txt};
\texttt{research/agents/codex/2026-11-11\_1203\_SAT\_POLYNOMIAL\_ASSUMPTION\_AUDIT.tex} (bridge (12), the $S_A$ system,
the "single next target");
\texttt{research/agents/codex/2026-11-11\_1331\_TRANSFER\_RECOVERY\_AND\_RATE\_AMPLIFICATION.tex}
(reflection eq.~(28)--(31), the Pich--Santhanam witnessing transfer);
LSPV-1/360 (TQBF-oracle nonrelativization of B1/B2/B3) and LSPV-1/361 (ER universal reflection,
FGPR bibliography entry); the repository search below; the prior Claude reports in this directory.
The newest LSPV commits (LSPV-1/610, LSPV-2/591, LSPV-3/606, LSPV-4/578) were surveyed by log only. They concern
native /341 checker semantics, and they neither bear on nor contradict the bridge question.

\paragraph{DUPLICATION / CONTRADICTION AUDIT.}
A repository search for \texttt{Grosser}, \texttt{Provable Reductions}, \texttt{TR26-109} and \texttt{TFNP}
found FGPR only as a bibliography item in LSPV-1/361 and in archived MAIN prewrites, with no theorem applied.
The oracle-collapse barrier for black-box bridges, which this run first planned to prove, is already in LSPV-1/360,
so it is \emph{not} claimed here. The fixed-$A$ polynomial equivalence $Q^*_A\in\mathcal G_q\Leftrightarrow W_E\in\mathcal G_q$ is
already in the Codex 1331 report, eq.~(30). Theorem~\ref{thm:wp}(b)$\Leftrightarrow$(c) restates it. The new parts are
only (a), the WrongProof/FGPR form, and (d), independence of the choice of decider.
No contradiction with the corpus was found.

\paragraph{WHAT CHANGED RELATIVE TO ROOT MAIN.}
Nothing in \texttt{main.tex} is edited. Suggested ledger entry for LSPV-1/LSPV-5: ``FGPR provable-reduction route
to $B(\EF)$: exact reformulation of the reflection condition; FAILED TRANSFER as a weaker bridge.''

\paragraph{ADOPTED / REJECTED / RETIRED PRIOR COMPONENTS.}
Adopted: the $S_A$ system (Codex 1203), the reflection bridge (LSPV-1/361, Codex 1331), and ER$\subseteq R_{\mathrm{mix341}}$ (LSPV-1/360).
Rejected as a new route: FGPR as a source of a weaker bridge. Nothing is retired.

\paragraph{STALE OR OVERSTATED CLAIMS FOUND.}
None are overstated. One citation needs checking: LSPV-1/361 cites FGPR as \texttt{arXiv:2606.27931}. This run
verified the ECCC record TR26-109 (dated 26--28 June 2026) but could not verify that arXiv identifier.

\section{Barrier-transfer and PQ ledger}
Obstacle 1: \emph{truth versus provability of correctness} of a hypothetical algorithm. The transfer domain is
bounded arithmetic and witnessing (FGPR; Pich--Santhanam 2023, ECCC TR23-199). Both are directly
transferable theorems, and both relocate the obstacle rather than remove it.
Obstacle 2: a lower bound for the same $R\ge_p\mathrm{ER}$, which is at least as hard as superpolynomial $\EF$ lower bounds and
has been open since Cook--Reckhow.
PQ item 6 (special soundness: two conflicting openings yield a forbidden witness) is structurally parallel to
$\WP$, where a satisfying assignment plus an accepted refutation yields an explicit verifier error.
\status{ANALOGY}: the cryptographic version gets its binding from a hardness assumption, which is not admissible here.
Relativization: the composition relativizes trivially as stated. It is a statement about provability, so it inherits LSPV-1/360:
any unconditional proof of (a) must use the code of $A$, not black-box access to it.

\section{Status table}
\begin{tabular}{p{0.56\linewidth}p{0.36\linewidth}}
Claim & Status\\\hline
$Q_A$ is a p-bounded Cook--Reckhow system under $\SAT\in\mathrm P$ & PROVED/EXACT\\
Theorem~\ref{thm:wp} (a)$\Leftrightarrow$(b)$\Leftrightarrow$(c)$\Leftrightarrow$(d) & PROVED/EXACT given (K1)--(K4)\\
Corollary for $R_{\mathrm{mix341}}$ & CONDITIONAL (ER-in-$R_{\mathrm{mix341}}$ overhead)\\
FGPR as a weaker bridge & FAILED TRANSFER\\
Unconditional $B(\EF)$ / $B(R_{\mathrm{mix341}})$ & OPEN\\
Superpolynomial $\overline W_{R_{\mathrm{mix341}}}$ & OPEN\\
P versus NP & OPEN
\end{tabular}

\section{Single next target}
Test the program-to-provable-program normalization on the strongest known partial case, before trying to prove it.
Is there a relativized or bounded-arithmetic model (for example, a $\mathrm{PV}$ model in the style of
Cook--Krajicek, \emph{Consequences of the provability of $\mathrm{NP}\subseteq\mathrm{P/poly}$}, JSL 2007)
in which SAT is decided by a polynomial-time algorithm while no SAT decider is $\EF$-provably correct?
Such a model would show that $B(\EF)$ is unprovable by any method formalizable in that setting.
That would be a firm negative result for the boxed bridge, and it fits the falsification-first instruction.

\section*{Computation}
None this run. All claims are short derivations from cited theorems. No Python was run, so the 30-second rule is not triggered.
No LaTeX engine is installed in this container, so the source is ASCII LaTeX and has not been compiled.

\section*{References}
N.~Fleming, S.~Grosser, T.~Pitassi, R.~Robere, \emph{Provable Reductions in TFNP}, ECCC TR26-109 (2026),
\url{https://eccc.weizmann.ac.il/report/2026/109/}; Thm.~5.2, Prop.~3.6, Prop.~3.8, Defs.~2.4, 3.5, 3.7, 5.1.
J.~Pich, R.~Santhanam, \emph{Towards $P\neq NP$ from Extended Frege lower bounds}, ECCC TR23-199,
\url{https://eccc.weizmann.ac.il/report/2023/199}.
J.~Krajicek, P.~Pudlak, Quantified propositional calculi and fragments of bounded arithmetic, \emph{Z.~Math.~Logik} 36 (1990).
S.~Cook, R.~Reckhow, The relative efficiency of propositional proof systems, \emph{JSL} 44 (1979).
S.~Cook, J.~Krajicek, Consequences of the provability of $NP\subseteq P/poly$, \emph{JSL} 72 (2007).

\end{document}
% END_OF_VERIFIED_REPORT::CLAUDE::WRONGPROOF-BRIDGE-0001
